%% This document created by Scientific Word (R) Version 2.5


\documentclass[12pt]{article}
%%%%%%%%%%%%%%%%%%%%%%%%%%%%%%%%%%%%%%%%%%%%%%%%%%%%%%%%%%%%%%%%%%%%%%%%%%%%%%%%%%%%%%%%%%%%%%%%%%%%%%%%%%%%%%%%%%%%%%%%%%%%%%
\usepackage{sw20jart}

%TCIDATA{TCIstyle=Article/ART4.LAT,jart,sw20jart}

%TCIDATA{Created=Tue Nov 18 18:18:44 2003}
%TCIDATA{LastRevised=Wed Nov 19 09:39:06 2003}

\input{tcilatex}

\begin{document}


\section{Misc. Calculations}

\subsection{Length of a field line}

If we define $s$ as the longitudinal distance along the magnetic field line
\begin{equation}
\mathbf{dl}=ds\widehat{b}
\end{equation}

The elementary distance in the toroidal direction is then
\begin{eqnarray}
Rd\phi  &=&\mathbf{dl}\cdot \widehat{\phi } \\
&=&ds\widehat{b}\cdot \widehat{\phi }
\end{eqnarray}
so that
\begin{equation}
d\phi =\frac{ds\widehat{b}\cdot \widehat{\phi }}{R}
\end{equation}

We have
\begin{equation}
\mathbf{B}=I\nabla \phi +\nabla \psi \times \nabla \phi 
\end{equation}
so that
\begin{eqnarray}
\widehat{b} &\equiv &\frac{\mathbf{B}}{B} \\
&=&\frac{I\left( \psi \right) \nabla \phi +\nabla \psi \times \nabla \phi }{B%
} \\
&=&\frac{I\left( \psi \right) \widehat{\phi }+\left| \nabla \psi \right| 
\widehat{b}_{p}}{BR}
\end{eqnarray}
and
\begin{equation}
d\phi =\frac{dsI\left( \psi \right) }{BR^{2}}
\end{equation}

The length of the field line corresponding to an entire toroidal period
therefore is (for an axysymmetric plasma)
\begin{equation}
s_{\text{tot}}\left( \psi \right) =\frac{2\pi BR^{2}}{I\left( \psi \right) }
\end{equation}

\subsection{Computation of $R\left( \psi ,s\right) $ from field data}

Assuming that the functions $I\left( \psi \right) $ and $\left| \nabla \psi
\right| \left( \psi ,s\right) $ are known, we have
\begin{eqnarray}
dR &=&\mathbf{dl}\cdot \widehat{R} \\
&=&ds\widehat{b}\cdot \widehat{R} \\
&=&ds\frac{\nabla \psi \times \widehat{\phi }\cdot \widehat{R}}{BR}
\end{eqnarray}
and therefore
\begin{equation}
dR=ds\frac{\nabla \psi \cdot \widehat{Z}}{BR}
\end{equation}
so that 
\begin{equation}
R\left( \psi ,s\right) =R\left( \psi ,0\right) +\int_{0}^{s}\frac{\nabla
\psi \cdot \widehat{Z}}{BR}ds
\end{equation}

\end{document}
